% SHARD-ID: AQM-59-FINITE-FIELD-CIRCLE-QUANTUM-SPACE
% SHARD-TITLE: The Finite-Field Circle and Its Complete Quantum Space
% SHARD-SUMMARY: Computes the affine circle \(x^2+y^2=1\) over \(\mathbb F_q\), separating odd and even characteristic.
% SHARD-SUMMARY: Derives rational point counts, cotangent fibres, and the nonreduced characteristic-two pathology.
% SHARD-SUMMARY: Describes the complete closed-point Weyl quantum space and its finite rational-point truncations.
% SHARD-KEYWORDS: finite field, circle, conic, cotangent, closed points, Weyl net, rational truncation, characteristic two
\section{The Finite-Field Circle and Its Complete Quantum Space}
\label{sec:finite-field-circle-quantum-space}
\subsection{Status and Source Anchors}
Let \(k=\mathbb F_q\) and
\[
  A_q=k[x,y]/(x^2+y^2-1),\qquad C_q=\Spec A_q.
\]
This shard uses Stacks algebra lines 34103--34120 for quotient differentials,
34310--34340 for polynomial differentials, and 34388--34524 for the de Rham
complex.  The closed-point finite-residue Weyl layer is convention (y) and
AQM-28--AQM-30.  The product-ring comparison for finite rational truncations
is AQM-40 and AQM-58.
\subsection{Odd Characteristic: Geometry and Differentials}
Assume \(q\) is odd.  In
\[
  S=k[x,y],\qquad f=x^2+y^2-1,\qquad A_q=S/(f),
\]
the quotient exact sequence gives
\[
  \Omega^1_{A_q/k}=
  (A_q\,dx\oplus A_q\,dy)/(2x\,dx+2y\,dy).
\]
Because \(2\in k^\times\), this is equivalently
\[
  x\,dx+y\,dy=0.
\]
At a \(k\)-rational point \(P=(a,b)\), the cotangent fibre is
\[
  T^*_{C_q/k,P}
  =
  (k\,dx\oplus k\,dy)/(a\,dx+b\,dy).
\]
Since \(a^2+b^2=1\), the pair \((a,b)\) is not \((0,0)\).  Thus the relation
has rank \(1\), and
\[
  \dim_k T^*_{C_q/k,P}=1,\qquad \dim_k T_{C_q/k,P}=1.
\]
Concrete bases are:
\[
\begin{array}{ll}
b\ne0: & T^*_{P}=k\,dx,\quad dy=-(a/b)\,dx,\\
b=0: & a=\pm1,\quad T^*_{P}=k\,dy,\quad dx=0.
\end{array}
\]
The same calculation over any closed point \(P\), with residue field
\(\kappa(P)\), gives a one-dimensional cotangent fibre over \(\kappa(P)\).
\subsection{Odd Characteristic: Rational Points as \texorpdfstring{\(q\)}{q} Varies}
The point \((1,0)\) lies on \(C_q\).  Draw a line of slope \(t\in k\) through
it:
\[
  y=t(x-1).
\]
Substitute into \(x^2+y^2=1\):
\[
  x^2+t^2(x-1)^2-1=(x-1)\bigl((1+t^2)x+(1-t^2)\bigr).
\]
One intersection is \(x=1\).  The second, when \(1+t^2\ne0\), is
\[
  x=\frac{t^2-1}{t^2+1},\qquad
  y=t(x-1)=\frac{-2t}{t^2+1}.
\]
This parametrizes all rational points except \((1,0)\), unless
\(1+t^2=0\), in which case the line passes through a missing point at
infinity and gives no second affine point.

Therefore
\[
  |C_q(k)|=1+\#\{t\in k:1+t^2\ne0\}.
\]
If \(-1\) is not a square in \(k\), then \(1+t^2\) never vanishes and
\[
  |C_q(k)|=q+1.
\]
If \(-1\) is a square, there are two roots of \(1+t^2\), and
\[
  |C_q(k)|=q-1.
\]
Equivalently, for odd \(q\),
\[
  |C_q(\mathbb F_q)|=q-\chi_q(-1),
\]
where \(\chi_q\) is the quadratic character on \(k^\times\).
The same calculation over the extension field \(\mathbb F_{q^m}\) gives
\[
  a_m(q):=|C_q(\mathbb F_{q^m})|=q^m-\chi_{q^m}(-1).
\]
If \(-1\) is already a square in \(k\), then \(\chi_{q^m}(-1)=1\) for every
\(m\), so \(a_m(q)=q^m-1\).  If \(-1\) is not a square in \(k\), then it
becomes a square exactly over even-degree extensions, so
\[
  a_m(q)=q^m-(-1)^m.
\]
\subsection{Even Characteristic: The Double Line}
Assume \(q\) is even.  Then
\[
  x^2+y^2-1=(x+y+1)^2.
\]
Thus
\[
  A_q=k[x,y]/((x+y+1)^2)
\]
is nonreduced.  Its reduced coordinate ring is
\[
  A_{q,\mathrm{red}}=k[x,y]/(x+y+1)\cong k[x],
\]
so the topological closed points are those of an affine line.  Rationally,
\[
  C_q(k)=\{(a,b)\in k^2:a+b+1=0\},
\]
and hence
\[
  |C_q(k)|=q.
\]
Over \(\mathbb F_{q^m}\), the same reduced line has
\[
  a_m(q)=q^m
\]
points.
The derivative pathology is sharper.  The relation is \(g^2=0\), where
\(g=x+y+1\), but
\[
  d(g^2)=2g\,dg=0
\]
in characteristic \(2\).  Hence the quotient differential relation imposes
nothing:
\[
  \Omega^1_{A_q/k}=A_q\,dx\oplus A_q\,dy.
\]
At a rational point \(P\), the residue field kills the nilpotent \(g\), but
the cotangent fibre remains
\[
  T^*_{C_q/k,P}=k\,dx\oplus k\,dy,
\]
of dimension \(2\).  The reduced line would have the relation \(dx+dy=0\) and
one-dimensional cotangent fibre.  The unreduced circle therefore has the same
closed-point Weyl sites as a line but a thicker cotangent geometry.
\subsection{Closed Points and the Complete Quantum Space}
The complete finite-residue quantum space is not only the rational-point set.
Let \(|C_q|_{\rm cl}\) be the closed points of \(C_q\).  For each closed point
\(P\), convention (y) assigns
\[
  E_P=\kappa(P)^2,\qquad
  \mathcal H_P=\ell^2(\kappa(P)),\qquad
  \mathcal A_P=\operatorname{End}_{\mathbb C}\mathcal H_P.
\]
The global algebraic quasi-local observable algebra is
\[
  \mathcal A(C_q)=
  \varinjlim_{S\Subset |C_q|_{\rm cl}}
  \bigotimes_{P\in S}\operatorname{End}_{\mathbb C}\ell^2(\kappa(P)).
\]
Equivalently, the Weyl label group is
\[
  E(C_q)=\bigoplus_{P\in |C_q|_{\rm cl}}\kappa(P)^2.
\]
After choosing the zero reference vector at every closed point, the
incomplete tensor Hilbert representation is
\[
  \mathcal H^0(C_q)=
  \ell^2\!\left(\bigoplus_{P\in |C_q|_{\rm cl}}\kappa(P)\right).
\]
This is the complete closed-point quantum space in the current report
convention.

Let \(N_d(q)\) be the number of degree-\(d\) closed points.  If
\[
  a_m(q)=|C_q(\mathbb F_{q^m})|,
\]
then every degree-\(d\) closed point contributes \(d\) points over
\(\mathbb F_{q^m}\) exactly when \(d\mid m\).  Hence
\[
  a_m(q)=\sum_{d\mid m}d\,N_d(q),
  \qquad
  N_d(q)=\frac1d\sum_{r\mid d}\mu(r)\,a_{d/r}(q).
\]
Thus the complete quantum space contains \(N_d(q)\) sites of local Hilbert
dimension \(q^d\) for each \(d\ge1\).

The rational truncation keeps only \(N_1(q)=|C_q(k)|\) sites:
\[
  \mathcal H_{\rm rat}(C_q)=
  \bigotimes_{P\in C_q(k)}\ell^2(k),
  \qquad
  \dim\mathcal H_{\rm rat}(C_q)=q^{|C_q(k)|}.
\]
For odd \(q\), this is \(q^{q-\chi_q(-1)}\).  For even \(q\), it is \(q^q\).
\subsection{Increasing \texorpdfstring{\(q\)}{q}: First Values}
\[
\begin{array}{c|c|c|c}
q & \operatorname{char} k & |C_q(k)| & \dim \mathcal H_{\rm rat}\\ \hline
2 & 2 & 2 & 2^2\\
3 & \text{odd},\ -1\text{ nonsquare} & 4 & 3^4\\
4 & 2 & 4 & 4^4\\
5 & \text{odd},\ -1\text{ square} & 4 & 5^4\\
7 & \text{odd},\ -1\text{ nonsquare} & 8 & 7^8\\
8 & 2 & 8 & 8^8\\
9 & \text{odd},\ -1\text{ square} & 8 & 9^8\\
11 & \text{odd},\ -1\text{ nonsquare} & 12 & 11^{12}\\
13 & \text{odd},\ -1\text{ square} & 12 & 13^{12}
\end{array}
\]
This table is only the rational finite truncation.  The complete quantum space
also has degree-\(d\) closed-point sites for every \(d\ge2\), counted by the
M\"obius formula above.
\subsection{Derivative Dynamics Versus Rational Truncation}
There is a subtle but important distinction.  If one replaces the rational
point set \(C_q(k)\) by its coordinate algebra of functions
\[
  k^{C_q(k)},
\]
then AQM-58 applies: this product ring has
\[
  \Omega^1_{k^{C_q(k)}/k}=0.
\]
So the rational finite set has trivial de Rham CSS dynamics.

The scheme \(C_q=\Spec A_q\) is different.  For odd \(q\), it is a smooth
affine curve with one-dimensional cotangent fibre at every closed point.  For
even \(q\), it is a nonreduced double line with two-dimensional cotangent
fibres before reduction.  Thus the complete closed-point Weyl quantum space
is built from residue fields at closed points, while the derivative geometry
is carried by \(\Omega^1_{A_q/k}\).  The former supplies the quantum sites;
the latter is the candidate source of dynamics.
